\chapter{Equações Diferenciais Lineares}\label{ch:b1:diffeq}

Encontradas pela primeira vez no volume do ensino médio, as equações
diferenciais são tratadas aqui com demonstrações completas e em maior
generalidade: equações lineares de primeira ordem com coeficientes
variáveis (resolvidas completamente pelo método da \omterm{thm:b1:diffeq:voc}{variação das
constantes}) e equações lineares de segunda ordem com coeficientes
constantes, o modelo das oscilações. Os dois casos exibem a mesma
estrutura: \emph{solução geral $=$ uma solução particular $+$ solução
geral da equação homogênea}.

\section{Equações lineares de primeira ordem}

\begin{definition}\label{def:b1:diffeq:linear1}
Sejam $I$ um intervalo e $a, b \colon I \to \R$ (ou $\C$)
funções contínuas. A equação
\[
(E)\colon\quad y' + a(x)\,y = b(x),
\]
na função incógnita derivável $y \colon I \to \R$ (ou $\C$),
é uma \emph{equação diferencial linear de primeira
ordem}\index{equação diferencial!linear de primeira ordem}. A
equação $(H)\colon y' + a(x) y = 0$ é a sua equação
\emph{homogênea}.
\end{definition}

\begin{theorem}[Resolvendo a equação homogênea]\label{thm:b1:diffeq:homogeneous1}
Seja $A$ uma primitiva de $a$ em $I$ (ela existe:
\cref{ch:b1:integration}). As soluções de $(H)$ em $I$ são exatamente
as funções
\[
y(x) = \lambda\, \eu^{-A(x)}, \qquad \lambda \in \R \text{ (ou} \C).
\]
\end{theorem}

\begin{proof}
Essas funções são soluções: $y' = -\lambda A' \eu^{-A} = -a y$.
Reciprocamente, seja $y$ uma solução de $(H)$ e ponha $z(x) = y(x)\, \eu^{A(x)}$.
Então
\[
z' = y' \eu^{A} + y\, a\, \eu^{A} = (y' + ay)\, \eu^{A} = 0 ,
\]
de modo que $z$ é constante no intervalo $I$, digamos $z = \lambda$: $y =
\lambda \eu^{-A}$. (Note a lógica: nenhuma solução se perde, porque
\emph{toda} solução foi escrita na forma anunciada.)
\end{proof}

\begin{example}[Uma equação homogênea com coeficiente variável]\label{ex:b1:diffeq:varcoeff}
Resolva $y' + (\cos x)\,y = 0$ em $\R$. Uma primitiva de $a(x) =
\cos x$ é $A(x) = \sin x$, de modo que as soluções são
\[
y(x) = \lambda\,\eu^{-\sin x}, \qquad \lambda \in \R .
\]
Duas leituras. Toda solução é periódica (de período $2\pi$) e nunca
se anula, a menos que $\lambda = 0$ --- o sinal de $\lambda$ é para
sempre o sinal de $y$, pois uma exponencial não pode cruzar o zero. E
a solução que passa por $y(0) = y_0$ é $y_0\eu^{-\sin x}$: exatamente
uma curva da família por cada ponto inicial, o retrato
unidimensional do \cref{thm:b1:diffeq:voc}~(2).
\end{example}

\begin{theorem}[Variação das constantes; problema de Cauchy]\label{thm:b1:diffeq:voc}
Com as notações acima:
\begin{enumerate}
  \item As soluções de $(E)$ em $I$ são exatamente
        \[
        y(x) = \Bigl(\lambda + \int_{x_0}^{x} b(t)\, \eu^{A(t)} \dd t
        \Bigr)\, \eu^{-A(x)},
        \qquad \lambda \in \R,
        \]
        em que $x_0 \in I$ está fixado. Equivalentemente: solução geral de
        $(H)$ mais uma solução particular de $(E)$.
  \item Para todos $x_0 \in I$ e $y_0$, o \emph{problema de
        Cauchy}\index{problema de Cauchy} ``$(E)$ e $y(x_0) = y_0$''
        tem exatamente uma solução em $I$.
\end{enumerate}
\end{theorem}

\begin{proof}
(1) Seguindo o método chamado \emph{variação das
constantes}\index{variação das constantes}, procuremos soluções da
forma $y = \mu(x)\, \eu^{-A(x)}$ com $\mu$ derivável --- nenhuma
generalidade se perde, pois toda função em $I$ pode ser assim escrita
($\mu = y\,\eu^{A}$). Substituindo,
\[
y' + ay = \mu' \eu^{-A} - \mu a \eu^{-A} + a\mu\eu^{-A}
= \mu'\, \eu^{-A},
\]
de modo que $y$ é solução de $(E)$ se, e somente se, $\mu'(x) = b(x)\,\eu^{A(x)}$, se
e somente se $\mu(x) = \lambda + \int_{x_0}^x b(t)\eu^{A(t)}\dd t$ para
alguma constante $\lambda$ (duas primitivas de uma mesma função
contínua num intervalo diferem por uma constante).

(2) Na fórmula, $y(x_0) = \lambda\,\eu^{-A(x_0)}$: a condição
$y(x_0) = y_0$ determina $\lambda = y_0 \eu^{A(x_0)}$ de modo único.
\end{proof}

\begin{example}\label{ex:b1:diffeq:order1}
Resolva $y' + \dfrac{y}{x} = x^2$ em $I = \intoo{0}{+\infty}$. Aqui
$a(x) = \frac 1x$, $A(x) = \ln x$, $\eu^{-A(x)} = \frac 1x$.
Soluções homogêneas: $\frac{\lambda}{x}$. \omterm{thm:b1:diffeq:voc}{Variação das constantes}:
$\mu'(x) = x^2 \cdot x = x^3$, de modo que $\mu = \frac{x^4}{4} + \lambda$, e
\[
y(x) = \frac{x^3}{4} + \frac{\lambda}{x}, \qquad \lambda \in \R .
\]
Com a condição inicial $y(1) = 0$: $\lambda = -\frac14$.
\emph{Verificação:} $y' + \frac yx = \frac{3x^2}{4} - \frac{\lambda}{x^2} +
\frac{x^2}{4} + \frac{\lambda}{x^2} = x^2$.
\end{example}

\begin{example}[Adivinhar bate integrar]\label{ex:b1:diffeq:guess}
Resolva $y' + 2x\,y = x$ em $\R$. A \omterm{thm:b1:diffeq:voc}{variação das constantes} funciona ($A =
x^2$, $\mu' = x\,\eu^{x^2}$, $\mu = \frac12\eu^{x^2} + \lambda$),
mas observar que a \emph{constante} $y_p = \frac12$ resolve a
equação ($0 + 2x\cdot\frac12 = x$) é mais rápido. Com as
soluções homogêneas $\lambda\,\eu^{-x^2}$:
\[
y(x) = \frac12 + \lambda\,\eu^{-x^2}, \qquad \lambda \in \R .
\]
Toda solução converge para $\frac12$ extremamente depressa quando $x \to
\pm\infty$: a solução particular constante é um
\emph{equilíbrio} ao qual todas as soluções se juntam. A ideia: antes
de lançar o método geral, gaste dez segundos procurando uma
solução particular óbvia (constante, monômio, múltiplo do
lado direito); o teorema de estrutura conclui então o trabalho.
\end{example}

\begin{remark}[Os intervalos importam]
O teorema vive num \emph{intervalo} em que $a$ e $b$ são
contínuas. Para $y' + \frac yx = 0$ em $\R^*$, as soluções são
$\frac{\lambda}{x}$ em $\intoo{0}{+\infty}$ e
$\frac{\mu}{x}$ em $\intoo{-\infty}{0}$, com constantes
\emph{independentes}: não há razão para uma única fórmula colar através
da singularidade em $0$.
\end{remark}

\begin{example}[Um lado direito complexo, duas respostas reais]\label{ex:b1:diffeq:complexrhs}
Resolva $y' - y = \cos x$ e $y' - y = \sin x$ de uma só vez. Trabalhe
em $\C$ com o lado direito $\eu^{\iu x}$: tentar $y_p =
c\,\eu^{\iu x}$ dá $c(\iu - 1)\eu^{\iu x} = \eu^{\iu x}$, de modo que
\[
c = \frac1{\iu - 1} = \frac{-1 - \iu}2,
\qquad
y_p = -\frac{(1 + \iu)(\cos x + \iu\sin x)}2
= \frac{\sin x - \cos x}2 + \iu\,\frac{-\sin x - \cos x}2 .
\]
Como a equação tem coeficientes reais, as partes real e imaginária
se separam: $\frac{\sin x - \cos x}2$ resolve $y' - y = \cos x$,
e $-\frac{\sin x + \cos x}2$ resolve $y' - y = \sin x$ (verifique
a primeira: a derivada $\frac{\cos x + \sin x}2$, menos a
função, dá $\cos x$). Uma linha complexa substituiu duas execuções da
\omterm{thm:b1:diffeq:voc}{variação das constantes} --- a mesma economia que o
\cref{met:b1:diffeq:particular} sistematiza para a segunda ordem,
e um dividendo recorrente do \cref{ch:b1:complex}.
\end{example}

\section{Equações lineares de segunda ordem com coeficientes constantes}

\begin{definition}\label{def:b1:diffeq:linear2}
Sejam $a, b \in \R$ e $f \colon I \to \R$ contínua. A equação
\[
(E)\colon\quad y'' + a\,y' + b\,y = f(x)
\]
é uma \emph{equação linear de segunda ordem com coeficientes
constantes}\index{equação diferencial!linear de segunda ordem};
$(H)\colon y'' + ay' + by = 0$ é a sua equação homogênea, e
$\chi(r) = r^2 + ar + b$, o seu \emph{polinômio
característico}\index{polinômio característico}.
\end{definition}

\begin{theorem}[Soluções homogêneas]\label{thm:b1:diffeq:homogeneous2}
Seja $\Delta = a^2 - 4b$ o discriminante de $\chi$. As soluções
reais de $(H)$ em $\R$ são:
\begin{enumerate}
  \item se $\Delta > 0$, com $r_1 \neq r_2$ as duas raízes reais:
        $\;y = \lambda\, \eu^{r_1 x} + \mu\, \eu^{r_2 x}$;
  \item se $\Delta = 0$, com $r_0$ a raiz dupla:
        $\;y = (\lambda + \mu x)\, \eu^{r_0 x}$;
  \item se $\Delta < 0$, com raízes $\alpha \pm \iu\omega$
        ($\omega > 0$):
        $\;y = \eu^{\alpha x} (\lambda \cos\omega x + \mu
        \sin\omega x)$;
\end{enumerate}
em cada caso com $(\lambda, \mu)$ percorrendo $\R^2$.
\end{theorem}

\begin{proof}
Observe primeiro que, para $r \in \C$, $x \mapsto \eu^{rx}$ é solução de $(H)$
se, e somente se, $\chi(r) = 0$ (substitua: $(r^2 + ar + b)\eu^{rx} =
0$). É por isso que as exponenciais são o primeiro palpite natural:
a derivação age sobre $\eu^{rx}$ como a multiplicação pelo número
$r$, de modo que a equação diferencial se torna a equação numérica
$\chi(r) = 0$ --- todo o problema analítico fica comprimido em
encontrar as raízes de uma equação do segundo grau.

O passo-chave é uma mudança de incógnita que reduz a ordem. Seja $r$
uma raiz (possivelmente complexa) de $\chi$ e escreva $y = z\, \eu^{rx}$,
o que não faz perder generalidade. Então
\[
y'' + ay' + by
= \bigl(z'' + (2r + a) z' + \chi(r) z\bigr)\eu^{rx}
= \bigl(z'' + (2r + a) z'\bigr)\eu^{rx},
\]
de modo que $(H)$ se torna a equação de \emph{primeira ordem} $u' + (2r + a) u = 0$
para $u = z'$.

\emph{Caso $\Delta \neq 0$:} escolha $r = r_1$; então $2r_1 + a = r_1 -
r_2$ (pois $r_1 + r_2 = -a$). Pelo
\cref{thm:b1:diffeq:homogeneous1}, $z' = c\,\eu^{(r_2 - r_1)x}$ para
alguma constante $c$; integrando em $\R$, $z = \mu\, \eu^{(r_2 - r_1)x}
+ \lambda$ com $\mu = \frac{c}{r_2 - r_1}$ e, portanto, $y = z\,
\eu^{r_1 x} = \lambda\,\eu^{r_1x} + \mu\,\eu^{r_2x}$.
Quando $\Delta < 0$, as raízes são $\alpha \pm \iu\omega$ e as
soluções complexas são $y = c_1\eu^{(\alpha+\iu\omega)x} +
c_2\eu^{(\alpha-\iu\omega)x}$ com $c_1, c_2 \in \C$. Quais delas assumem valores reais? Como $\conj{\eu^{(\alpha+\iu\omega)x}} =
\eu^{(\alpha-\iu\omega)x}$, o conjugado de $y$ é $\conj{c_2}\,
\eu^{(\alpha+\iu\omega)x} + \conj{c_1}\,\eu^{(\alpha-\iu\omega)x}$,
e $y = \conj y$ para todo $x$ força $c_2 = \conj{c_1}$ (as duas
exponenciais são linearmente independentes: avalie em dois pontos, ou
compare em $x=0$ depois de dividir por $\eu^{\alpha x}$). Escrevendo $c_1
= \frac{\lambda - \iu\mu}2$ com $\lambda, \mu$ reais:
\[
y = 2\,\Re\Bigl(c_1\,\eu^{\alpha x}(\cos\omega x +
\iu\sin\omega x)\Bigr)
= \eu^{\alpha x}\bigl(\lambda\cos\omega x + \mu\sin\omega
x\bigr),
\]
e, reciprocamente, toda função desse tipo é solução (parte real de uma
solução complexa de uma equação real): o espaço das soluções reais é
o anunciado.

\emph{Caso $\Delta = 0$:} $r = r_0$, $2r_0 + a = 0$, de modo que $z'' = 0$:
$z = \lambda + \mu x$ e $y = (\lambda + \mu x)\eu^{r_0 x}$.
\end{proof}

\begin{example}[Um problema de Cauchy, do início ao fim]\label{ex:b1:diffeq:cauchy2}
Resolva $y'' - 3y' + 2y = 0$ com $y(0) = 0$, $y'(0) = 1$. O
\omterm{def:b1:diffeq:linear2}{polinômio característico} $r^2 - 3r + 2 = (r - 1)(r - 2)$ tem as
raízes reais $1$ e $2$: solução geral $y = \lambda\eu^{x} +
\mu\eu^{2x}$. As duas condições dão o sistema linear
\[
\lambda + \mu = 0, \qquad \lambda + 2\mu = 1 ,
\]
de modo que $\mu = 1$, $\lambda = -1$:
\[
y(x) = \eu^{2x} - \eu^{x} .
\]
Verificação: $y(0) = 0$; $y' = 2\eu^{2x} - \eu^x$ tem $y'(0) = 1$; e
$y'' - 3y' + 2y = (4 - 6 + 2)\eu^{2x} + (-1 + 3 - 2)\eu^x = 0$.
Note a forma da resposta: perto de $-\infty$ domina o modo lento
$-\eu^x$; perto de $+\infty$, o modo rápido $\eu^{2x}$.
Ler as soluções como superposições de modos com taxas diferentes de
decaimento ou crescimento é o hábito rentável --- é assim que a
separação transitório/regime permanente do problema de fim de semana é
organizada.
\end{example}

\begin{omfigure}
\begin{tikzpicture}[xscale=1.35, yscale=1.55]
  \draw[->, thin] (0,0) -- (6.4,0) node[below] {$x$};
  \draw[->, thin] (0,-0.85) -- (0,1.25) node[left] {$y$};
  \draw[omDef, thick, domain=0:6.2, samples=400, smooth]
    plot (\x, {exp(-0.35*\x)*cos(4*\x r)});
  \draw[omThm, thick, domain=0:6.2, samples=120, smooth]
    plot (\x, {(1+1.8*\x)*exp(-1.1*\x)});
  \draw[omProp, thick, domain=0:6.2, samples=120, smooth]
    plot (\x, {1.15*exp(-0.45*\x) - 0.15*exp(-2.5*\x)});
  \node[omDef, font=\small] at (2.1,-0.62)
    {$\Delta < 0$: oscila};
  \node[omThm, font=\small] at (0.95,1.1) {$\Delta = 0$};
  \node[omProp, font=\small] at (4.9,0.38)
    {$\Delta > 0$: sem oscilação};
\end{tikzpicture}

{\small Os três regimes de $y'' + ay' + by = 0$ com soluções
que decaem: oscilação amortecida (raízes complexas), retorno crítico
(raiz dupla), decaimento superamortecido (duas raízes reais). Qual
regime ocorre lê-se apenas no sinal de $\Delta = a^2 - 4b$ ---
antes de resolver o que quer que seja.}
\end{omfigure}

\begin{theorem}[Estrutura e problema de Cauchy]\label{thm:b1:diffeq:structure2}
\begin{enumerate}
  \item Se $y_p$ é uma solução particular de $(E)$, as soluções de
        $(E)$ são exatamente $y_p + y_h$, com $y_h$ percorrendo as
        soluções de $(H)$.
  \item (Superposição) Se $y_1$ resolve $y'' + ay' + by = f_1$ e
        $y_2$ resolve $y'' + ay' + by = f_2$, então $y_1 + y_2$ resolve
        a equação com lado direito $f_1 + f_2$.
  \item Para todos $x_0 \in I$ e $(y_0, y_0')$, o \omterm{thm:b1:diffeq:voc}{problema de Cauchy}
        ``$(E)$, $y(x_0) = y_0$, $y'(x_0) = y_0'$'' tem exatamente uma
        solução em $I$. \emph{(Existência dada uma solução
        particular; unicidade em geral.)}
\end{enumerate}
\end{theorem}

\begin{proof}
(1) $y$ resolve $(E)$ se, e somente se, $y - y_p$ resolve $(H)$, pela
linearidade de $y
\mapsto y'' + ay' + by$. (2) é a mesma linearidade.

(3) Por (1), basta demonstrar que as constantes $(\lambda,
\mu)$ podem
sempre ser ajustadas, de modo único, a quaisquer dados $(y_0, y_0')$.
Transladando a variável, suponha $x_0 = 0$. No caso (1) do
\cref{thm:b1:diffeq:homogeneous2}, $y = \lambda\eu^{r_1x} +
\mu\eu^{r_2x}$ dá
\[
y(0) = \lambda + \mu, \qquad y'(0) = r_1\lambda + r_2\mu :
\]
um sistema linear em $(\lambda, \mu)$ cujo determinante é $r_2 -
r_1 \neq 0$; resolvendo-o explicitamente, $\mu = \frac{y_0' -
r_1y_0}{r_2 - r_1}$ e $\lambda = y_0 - \mu$: exatamente uma
solução. No caso (2), $y(0) = \lambda$ e $y'(0) = r_0\lambda +
\mu$: o sistema é triangular, com determinante $1$, resolvido por
$\lambda = y_0$, $\mu = y_0' - r_0y_0$. No caso (3), $y(0) =
\lambda$ e $y'(0) = \alpha\lambda + \omega\mu$: determinante
$\omega \neq 0$, resolvido por $\lambda = y_0$, $\mu = \frac{y_0' -
\alpha y_0}\omega$. Em cada caso, a \omterm{def:b1:logic:map}{aplicação} $(\lambda, \mu) \mapsto
(y(x_0), y'(x_0))$ é uma bijeção linear --- a linguagem do
\cref{ch:b1:linmaps} comprimirá esta verificação por casos numa única
frase.
\end{proof}

\begin{method}[Solução particular para $f(x) = P(x)\,\eu^{\gamma x}$]\label{met:b1:diffeq:particular}
Quando o lado direito é $P(x)\,\eu^{\gamma x}$ com $P$ um
polinômio e $\gamma \in \R$ (isso cobre polinômios, exponenciais e,
via $\gamma$ complexo ou superposição, $\cos$ e $\sin$), procure
uma solução particular da forma
\[
y_p(x) = x^{m}\, Q(x)\, \eu^{\gamma x},
\qquad
m = \text{multiplicidade de} \gamma \text{ como raiz de} \chi
\ (m = 0, 1 \text{ ou} 2),
\]
com $Q$ um polinômio do mesmo grau que $P$, cujos coeficientes
se encontram por substituição e identificação. Para $f = K\cos\omega x$
(ou $\sin$), resolva com lado direito $K\eu^{\iu\omega x}$ e tome
a parte real (resp.\ imaginária).
\end{method}

\begin{example}[Superposição em ação]\label{ex:b1:diffeq:superposition}
Resolva $y'' - y = \eu^{x} + 4$ em $\R$. Homogênea: $\chi(r) = r^2
- 1$, raízes $\pm1$, de modo que $y_h = \lambda\eu^x + \mu\eu^{-x}$. Separe
o lado direito e trate cada parcela pelo quadro do método.
\emph{Parcela $\eu^x$:} aqui $\gamma = 1$ é raiz simples de
$\chi$, então tente $y_1 = c\,x\,\eu^x$: nesse caso $y_1'' - y_1 = c(x +
2)\eu^x - cx\eu^x = 2c\,\eu^x$, o que dá $c = \frac12$.
\emph{Parcela $4$:} $\gamma = 0$ não é raiz; a constante $y_2 =
-4$ serve. Por superposição
(\cref{thm:b1:diffeq:structure2}~(2)):
\[
y = \frac{x\,\eu^x}2 - 4 + \lambda\,\eu^x + \mu\,\eu^{-x},
\qquad (\lambda, \mu) \in \R^2 .
\]
Note como as duas parcelas exigiram formas \emph{diferentes} ($m = 1$
contra $m = 0$): o teste de multiplicidade é aplicado a cada
expoente separadamente, o que é toda a razão de separar o
lado direito antes de adivinhar.
\end{example}

\begin{example}[A regra da multiplicidade em ação]\label{ex:b1:diffeq:multiplicity}
Resolva $y'' + y' = x$ em $\R$. O lado direito é $P(x)\eu^{0
\cdot x}$
com $P(x) = x$, e $\gamma = 0$ é raiz \emph{simples}
de $\chi(r) = r^2 + r = r(r + 1)$: logo $m = 1$, e o palpite correto
é $y_p = x\,(\alpha x + \beta) = \alpha x^2 + \beta x$, um grau
acima de $P$. Substituindo:
\[
y_p'' + y_p' = 2\alpha + (2\alpha x + \beta)
= 2\alpha x + (2\alpha + \beta) ,
\]
e a identificação com $x$ dá $\alpha = \frac12$, $\beta = -1$:
$y_p = \frac{x^2}2 - x$. Solução geral: $y = \frac{x^2}2 - x +
\lambda + \mu\,\eu^{-x}$. Se tivéssemos adivinhado $y_p = \alpha x + \beta$
(ignorando a multiplicidade), a substituição daria $y_p'' + y_p'
= \alpha$, uma constante --- nenhuma escolha de $\alpha, \beta$ pode igualar
$x$, e a falha é estrutural: as constantes já resolvem a equação
homogênea, de modo que são invisíveis ao lado esquerdo.
O fator $x^m$ existe precisamente para sair do espaço das soluções
homogêneas.
\end{example}

\begin{remark}[A apólice de seguro de trinta segundos]\label{rem:b1:diffeq:check}
Toda equação resolvida neste capítulo termina com uma verificação por
substituição, e isso não é decorativo. Um cálculo de equação
diferencial encadeia muitos passos pequenos (uma primitiva, uma regra
do produto, duas constantes), e um único erro de sinal se propaga
invisivelmente; substituir a fórmula final na equação captura
essencialmente todos eles ao custo de uma derivação.
Cultive o reflexo em três camadas: verifique a \emph{solução
particular} sozinha (a parte homogênea se cancela de qualquer modo),
verifique as \emph{condições iniciais} na solução completa e, quando
houver um parâmetro, verifique um \emph{valor degenerado} (a fórmula
para $\Omega$ geral reproduz a resposta conhecida em
$\Omega = 0$?). O hábito custa meio minuto e converte
``provavelmente certo'' em ``verificado''.
\end{remark}

\begin{remark}[Armadilhas frequentes]\label{rem:b1:diffeq:pitfalls}
\begin{enumerate}
  \item \emph{Normalize primeiro.} As fórmulas supõem que a equação
        se lê $y' + a(x)y = b(x)$ --- coeficiente $1$ em $y'$.
        Para $xy' - 2y = x^3$, divida por $x$ (num intervalo
        que evite $0$) antes de identificar $a$ e $b$, como no
        \cref{exo:b1:diffeq:2}.
  \item \emph{Uma constante por dimensão, fixada no fim.} A
        solução geral de primeira ordem carrega uma constante; a de
        segunda ordem, duas; as condições iniciais são impostas à
        solução \emph{completa} $y_p + y_h$, nunca a $y_h$
        sozinha --- impô-las antes de somar $y_p$ é o erro
        estrutural mais frequente.
  \item \emph{Atenção à multiplicidade.} Um palpite de solução
        particular que resolve a equação homogênea é invisível ao
        lado esquerdo; o fator $x^m$ do
        \cref{met:b1:diffeq:particular} não é opcional
        (\cref{ex:b1:diffeq:multiplicity}).
  \item \emph{Os intervalos fazem parte da resposta.} As soluções
        vivem em intervalos nos quais os coeficientes são contínuos;
        colar através de uma singularidade pode criar constantes
        espúrias (\cref{exo:b1:diffeq:12}) ou destruir a unicidade. ``Resolva
        em $\R^*$'' significa dois problemas independentes.
\end{enumerate}
\end{remark}

\begin{example}[Forçamento fora de ressonância]\label{ex:b1:diffeq:offresonance}
Resolva $y'' + 4y = \sin x$ em $\R$. Frequência natural $2$, frequência
de forçamento $1$: como $\iu$ \emph{não} é raiz de $\chi(r) = r^2
+ 4$, a multiplicidade é $m = 0$ e uma senoide simples basta.
Tentando $y_p = \alpha\sin x$ (nenhum cosseno é necessário: a equação não
tem termo em $y'$, e $\sin$ regenera $\sin$):
\[
y_p'' + 4y_p = -\alpha\sin x + 4\alpha\sin x = 3\alpha\sin x ,
\]
de modo que $\alpha = \frac13$, e a solução geral é
\[
y = \frac{\sin x}3 + \lambda\cos 2x + \mu\sin 2x .
\]
Toda solução permanece limitada: uma superposição de duas oscilações
nas frequências $1$ (forçada) e $2$ (natural). Compare com
o exemplo seguinte, em que forçar \emph{na} frequência natural
muda a própria forma da resposta.
\end{example}

\begin{example}[Uma oscilação forçada]\label{ex:b1:diffeq:oscillation}
Resolva $y'' + y = \cos x$, $y(0) = 0$, $y'(0) = 0$.

\emph{Homogênea:} $\chi(r) = r^2 + 1$, raízes $\pm\iu$: $y_h =
\lambda\cos x + \mu \sin x$.

\emph{Particular:} lado direito $\Re(\eu^{\iu x})$ com $\gamma =
\iu$ raiz simples de $\chi$: tente $z_p = c\, x\, \eu^{\iu x}$
($c \in \C$). Então $z_p'' + z_p = c\,(2\iu)\eu^{\iu x}$, o que é igual a
$\eu^{\iu x}$ para $c = \frac{1}{2\iu} = -\frac\iu2$. Logo, $z_p =
-\frac{\iu}{2} x (\cos x + \iu \sin x)$ e $y_p = \Re(z_p) =
\frac{x \sin x}{2}$.

\emph{Solução geral:} $y = \frac{x\sin x}{2} + \lambda\cos x +
\mu\sin x$. Condições: $y(0) = \lambda = 0$; $y' = \frac{\sin x +
x\cos x}{2} + \mu\cos x$, de modo que $y'(0) = \mu = 0$. Resposta: $y =
\frac{x\sin x}{2}$ --- uma oscilação cuja amplitude cresce linearmente:
o fenômeno da \emph{ressonância}\index{ressonância}, causado por forçar
o sistema na sua frequência natural.
\end{example}

\begin{omfigure}
\begin{tikzpicture}
\begin{axis}[omaxis, width=11cm, height=5.2cm,
    xmin=0, xmax=25, ymin=-13, ymax=13,
    xtick={0,5,10,15,20,25}, ytick=\empty]
  \addplot[omThm, thick, domain=0:25, samples=400] {x*sin(deg(x))/2};
  \addplot[dashed, gray, domain=0:25] {x/2};
  \addplot[dashed, gray, domain=0:25] {-x/2};
\end{axis}
\end{tikzpicture}

{\small \omterm{ex:b1:diffeq:oscillation}{Ressonância}: a solução $y = \frac{x \sin x}{2}$ de $y'' + y =
\cos x$
oscila entre as retas $y = \pm\frac x2$ (tracejadas), com
amplitude sempre crescente.}
\end{omfigure}

\begin{remark}[Interlúdio: a linearidade é uma geometria]\label{rem:b1:diffeq:interlude}
Olhe de novo a forma de cada \omterm{def:b1:logic:sets}{conjunto} solução deste capítulo: uma
solução especial mais um espaço de soluções homogêneas com uma
constante livre (primeira ordem) ou duas (segunda ordem). Os capítulos
de álgebra linear (\cref{ch:b1:vspaces,ch:b1:findim,ch:b1:linmaps})
fornecerão o vocabulário exato: a \omterm{def:b1:logic:map}{aplicação} $L(y) = y'' + ay' + by$ é
\emph{linear}, as suas soluções homogêneas formam o \emph{núcleo} de
$L$, um espaço vetorial cuja \emph{dimensão} é igual à ordem da equação
--- eis o conteúdo honesto de ``uma constante por ordem'' --- e
o \omterm{def:b1:logic:sets}{conjunto} solução de $L(y) = f$ é um \emph{subespaço afim}, um
transladado do núcleo. Até a \omterm{def:b1:logic:map}{aplicação} de Cauchy $(\lambda, \mu)
\mapsto (y(x_0), y'(x_0))$ do \cref{thm:b1:diffeq:structure2} é uma
bijeção linear entre dois planos, isto é, um sistema $2
\times 2$
invertível (\cref{ch:b1:matrices}). Nada neste capítulo
precisará ser refeito --- apenas rebatizado, e o rebatismo é o melhor
aquecimento possível para a álgebra linear:
toda definição abstrata de lá já ganhou o seu sustento aqui.
\end{remark}

\begin{remark}[Onde este capítulo é usado]\label{rem:b1:diffeq:whereused}
O teorema de estrutura --- as soluções de $(E)$ formam ``uma solução
particular mais as soluções de $(H)$'' --- é a primeira aparição
de um padrão que o \cref{ch:b1:vspaces,ch:b1:linmaps} nomeará: o
\omterm{def:b1:logic:sets}{conjunto} solução de $(H)$ é o \emph{núcleo} da \omterm{def:b1:logic:map}{aplicação} linear $y
\mapsto y'' + ay' + by$, e o \omterm{def:b1:logic:sets}{conjunto} solução de $(E)$ é um
transladado afim dele. O \omterm{def:b1:diffeq:linear2}{polinômio característico} reaparece como
o \omterm{def:b1:diffeq:linear2}{polinômio característico} de uma matriz no \cref{ch:b1:matrices}:
uma equação de segunda ordem é um sistema de primeira ordem $2 \times 2$ disfarçado, ponto de vista que o volume do segundo ano de graduação
sistematiza. As integrais exigidas pela \omterm{thm:b1:diffeq:voc}{variação das constantes} são
fornecidas pelo \cref{ch:b1:integration}, e o problema de fim de semana abaixo ---
o oscilador amortecido e forçado --- é o caso modelo de toda
questão de oscilação nas ciências, dos circuitos às pontes
suspensas.
\end{remark}

\section{Exercícios}

\begin{exercise}[$\star$]\label{exo:b1:diffeq:1}
Resolva em $\R$: $\;y' + 2y = \eu^{3x}$; depois o \omterm{thm:b1:diffeq:voc}{problema de Cauchy} $y(0) =
1$.
\end{exercise}

\begin{exercise}[$\star$]\label{exo:b1:diffeq:2}
Resolva em $\intoo{0}{+\infty}$: $\;x y' - 2y = x^3$ \emph{(coloque a
equação na forma normalizada primeiro)}.
\end{exercise}

\begin{exercise}[$\star$]\label{exo:b1:diffeq:3}
Resolva em $\R$, dando a solução geral real:
$\;y'' - 3y' + 2y = 0$; $\;y'' + 4y' + 4y = 0$; $\;y'' - 2y' + 5y =
0$.
\end{exercise}

\begin{exercise}[$\star$]\label{exo:b1:diffeq:4}
Resolva $y'' - y = x^2$ em $\R$ e depois o \omterm{thm:b1:diffeq:voc}{problema de Cauchy} $y(0) = 0$, $y'(0) = 1$.
\end{exercise}

\begin{exercise}[$\star\star$]\label{exo:b1:diffeq:5}
Resolva em $\intoo{-\frac\pi2}{\frac\pi2}$:
$\;y' + y\tan x = \sin 2x$.
\end{exercise}

\begin{exercise}[$\star\star$]\label{exo:b1:diffeq:6}
Resolva $y'' - 4y' + 3y = (2x + 1)\,\eu^{x}$ em $\R$. \emph{(Atenção à
multiplicidade: $1$ é raiz do \omterm{def:b1:diffeq:linear2}{polinômio característico}?)}
\end{exercise}

\begin{exercise}[$\star\star$]\label{exo:b1:diffeq:7}
Resolva $y'' + 4y = \sin 2x + x$ em $\R$ \emph{(superposição; trate
cada lado direito separadamente)}.
\end{exercise}

\begin{exercise}[$\star\star$]\label{exo:b1:diffeq:8}
Uma xícara de café à temperatura $T_0 = 80\,^\circ$C está numa sala a
$20\,^\circ$C. A lei do resfriamento de Newton diz que $T' = -k\,(T - 20)$ com
$k > 0$. Resolva para $T(t)$ e, sabendo que o café está a
$50\,^\circ$C após $10$ minutos, determine quando ele atinge
$25\,^\circ$C.
\end{exercise}

\begin{exercise}[$\star\star\star$]\label{exo:b1:diffeq:9}
(Oscilador amortecido) Para $\varepsilon \geq 0$, considere
$y'' + 2\varepsilon y' + y = 0$.
\begin{enumerate}
  \item Resolva para $\varepsilon \in \intco{0}{1}$, $\varepsilon = 1$
        e $\varepsilon > 1$.
  \item Mostre que, para $\varepsilon > 0$, toda solução tende a $0$
        em $+\infty$, e que, para $\varepsilon = 0$, as soluções não
        nulas não tendem.
  \item Para $\varepsilon \in \intoo{0}{1}$, mostre que os zeros de uma
        solução não nula são regularmente espaçados, com passo
        $\frac{\pi}{\sqrt{1 - \varepsilon^2}}$.
\end{enumerate}
\end{exercise}

\begin{exercise}[$\star\star\star$]\label{exo:b1:diffeq:10}
Encontre todas as funções $f \colon \R \to \R$, duas vezes deriváveis, tais
que
\[
\forall x, y \in \R, \qquad f(x + y) + f(x - y) = 2 f(x) f(y),
\]
com $f(0) \ne 0$ e $f$ não constante.
\emph{Sugestão: fixe $y$, derive duas vezes em relação a $x$ em $0$;
mostre que $f(0) = 1$ e $f'' = c f$ para alguma constante $c$; depois
resolva conforme o sinal de $c$ e verifique quais soluções satisfazem a
equação funcional.}
\end{exercise}

\begin{exercise}[$\star\star$]\label{exo:b1:diffeq:11}
(Equação de Euler) Resolva $x^2 y'' - x y' + y = 0$ em
$\intoo{0}{+\infty}$. \emph{Sugestão: ponha $z(t) = y(\eu^t)$, isto é,
substitua $x = \eu^t$, e mostre que $z$ satisfaz uma equação
linear com coeficientes constantes.}
\end{exercise}

\begin{exercise}[$\star\star\star$]\label{exo:b1:diffeq:12}
Considere a equação $x\,y' = 2y$ em toda a reta real, na
função incógnita derivável $y \colon \R \to \R$.
\begin{enumerate}
  \item Resolva em $\intoo{0}{+\infty}$ e em $\intoo{-\infty}{0}$.
  \item Mostre que, para \emph{quaisquer} constantes $a, b \in \R$, a
        função igual a $ax^2$ para $x \geq 0$ e a $bx^2$ para
        $x < 0$ é derivável em $\R$ e resolve a equação
        em toda parte.
  \item Conclua que o \omterm{def:b1:logic:sets}{conjunto} solução em $\R$ é uma família a dois
        parâmetros e explique por que isso não contradiz a
        unicidade do \cref{thm:b1:diffeq:voc}.
\end{enumerate}
\end{exercise}

\section{Problema: O oscilador amortecido e forçado}

\begin{problem}\label{pb:b1:diffeq:1}
Uma única equação governa uma massa presa a uma mola num meio
viscoso, a carga num circuito RLC e um prédio balançando ao vento:
\[
(E_\Omega)\colon\quad
x'' + 2\lambda x' + \omega_0^2\,x = A\cos(\Omega t),
\]
com $\lambda \geq 0$ o amortecimento, $\omega_0 > 0$ a frequência
natural, e $A > 0$, $\Omega > 0$ a amplitude e a frequência do
forçamento. Este problema extrai o seu comportamento completo: o
decaimento dos transitórios, o único regime permanente periódico, a
curva de \omterm{ex:b1:diffeq:oscillation}{ressonância} e a sua nitidez (o \emph{\omterm{pb:b1:diffeq:1}{fator de qualidade}}), os
batimentos do caso não amortecido e o balanço de energia que sustenta a
oscilação.\index{ressonância}\index{oscilador}
Salvo menção em contrário, $0 < \lambda < \omega_0$ (regime
subamortecido) e escrevemos $\omega_d = \sqrt{\omega_0^2 - \lambda^2}$.

\medskip
\noindent\textbf{Parte I --- O oscilador livre.} Aqui $A = 0$.
\begin{enumerate}
  \item Resolva a equação homogênea $(H)$ para $0 < \lambda <
        \omega_0$ e para $\lambda = 0$. (Os regimes $\lambda
        \geq \omega_0$ foram tratados no \cref{exo:b1:diffeq:9};
        cite-os.)
  \item Mostre que, para todo $\lambda > 0$, todas as soluções de $(H)$
        tendem a $0$ em $+\infty$ --- nos três regimes.
  \item Defina a energia $\mathcal E(t) = \frac12 x'(t)^2 +
        \frac12\omega_0^2\,x(t)^2$ ao longo de uma solução de $(H)$. Mostre que
        $\mathcal E'(t) = -2\lambda\,x'(t)^2 \leq 0$ e deduza
        (sem resolver nada) que o \omterm{thm:b1:diffeq:voc}{problema de Cauchy}
        ``$(H)$, $x(t_0) = x'(t_0) = 0$'' só tem a solução
        nula, para todo $\lambda \geq 0$.
  \item Para $0 < \lambda < \omega_0$, escreva a solução não nula como
        $x(t) = R\,\eu^{-\lambda t}\cos(\omega_d t - \varphi)$ e
        seja $T_d = \frac{2\pi}{\omega_d}$ o pseudoperíodo.
        Mostre que $x(t + T_d) = \eu^{-\lambda T_d}\,x(t)$: cada
        oscilação é a anterior encolhida pelo fator constante
        $\eu^{-\delta}$, $\delta = \frac{2\pi\lambda}{\omega_d}$
        (o \emph{decremento logarítmico}). Calcule $\delta$ para
        $\omega_0 = 1$, $\lambda = 0.1$.
  \item Defina o \emph{\omterm{pb:b1:diffeq:1}{fator de qualidade}} $Q =
        \dfrac{\omega_0}{2\lambda}$\index{fator de qualidade}. Mostre
        que, após o tempo $\frac1\lambda$ (uma multiplicação da
        amplitude por $1/\eu$), o oscilador completou
        $\frac{\omega_d}{2\pi\lambda}$ pseudoperíodos, o que, para
        amortecimento fraco ($\lambda \ll \omega_0$), é aproximadamente
        $\frac Q\pi$: o \omterm{pb:b1:diffeq:1}{fator de qualidade} conta, a menos de $\pi$, as
        oscilações sobreviventes antes de a amplitude decair por um
        fator $\eu$.
\end{enumerate}

\medskip
\noindent\textbf{Parte II --- O regime permanente.} Agora $A > 0$ e
$\lambda > 0$.
\begin{enumerate}
  \item Procure uma solução particular como parte real de
        $z\,\eu^{\iu\Omega t}$ com $z \in \C$
        (\cref{met:b1:diffeq:particular}). Mostre que isso funciona
        com
        \[
        z = \frac{A}{\omega_0^2 - \Omega^2 + 2\iu\lambda\Omega} .
        \]
  \item Deduza o regime permanente na forma amplitude--fase:
        $x_p(t) = R(\Omega)\cos\bigl(\Omega t -
        \varphi(\Omega)\bigr)$ com
        \[
        R(\Omega) = \frac{A}{\sqrt{(\omega_0^2 - \Omega^2)^2 +
        4\lambda^2\Omega^2}},
        \qquad
        \tan\varphi = \frac{2\lambda\Omega}{\omega_0^2 -
        \Omega^2},
        \quad \varphi \in \intoo0\pi .
        \]
  \item Interprete os dois regimes extremos: calcule os limites de
        $R$ e $\varphi$ quando $\Omega \to 0^+$ (resposta quase
        estática $A/\omega_0^2$, fase $0$) e quando $\Omega \to
        +\infty$ ($R \sim A/\Omega^2 \to 0$, fase $\to \pi$: a
        massa se move em oposição a um forçamento rápido demais).
  \item Mostre que \emph{toda} solução de $(E_\Omega)$ é $x_p$
        mais uma solução de $(H)$ e, portanto, converge para o
        regime permanente $x_p$ quando $t \to +\infty$, quaisquer que
        sejam as condições iniciais: depois que o transitório morre,
        o oscilador não guarda memória de como começou.
  \item Mostre que $x_p$ é a \emph{única} solução periódica de
        $(E_\Omega)$.
  \item Leve um \omterm{thm:b1:diffeq:voc}{problema de Cauchy} até o fim: para $x'' + 2x' + 2x =
        \cos t$ com $x(0) = x'(0) = 0$, mostre que a solução é
        \[
        x(t) = \frac{\cos t + 2\sin t}5
        - \eu^{-t}\,\frac{\cos t + 3\sin t}5 ,
        \]
        e identifique as partes transitória e permanente.
\end{enumerate}

\medskip
\noindent\textbf{Parte III --- A curva de \omterm{ex:b1:diffeq:oscillation}{ressonância}.} Estudo de
$\Omega \mapsto R(\Omega)$ em $\intoo0{+\infty}$.
\begin{enumerate}
  \item Pondo $u = \Omega^2$ e $g(u) = (\omega_0^2 - u)^2 +
        4\lambda^2 u$, mostre que: se $2\lambda^2 < \omega_0^2$, então $R$
        atinge um máximo estrito na \emph{frequência de \omterm{ex:b1:diffeq:oscillation}{ressonância}}
        $\Omega_r = \sqrt{\omega_0^2 - 2\lambda^2}$, com
        \[
        R_{\max} = R(\Omega_r)
        = \frac{A}{2\lambda\sqrt{\omega_0^2 - \lambda^2}} .
        \]
  \item Mostre que $\dfrac{R_{\max}}{R(0)} = Q\,\bigl(1 -
        \tfrac{\lambda^2}{\omega_0^2}\bigr)^{-1/2} \geq Q$: na
        \omterm{ex:b1:diffeq:oscillation}{ressonância}, o forçamento é amplificado (essencialmente) pelo
        \omterm{pb:b1:diffeq:1}{fator de qualidade}.
  \item Demonstre que a amplitude da \emph{velocidade} $V(\Omega) =
        \Omega\,R(\Omega)$ é máxima exatamente em $\Omega =
        \omega_0$ (e não em $\Omega_r$), e que ali a fase vale
        $\varphi(\omega_0) = \frac\pi2$: em $\Omega = \omega_0$
        a velocidade está exatamente em fase com a força.
  \item (Largura de banda) Resolva $g(u) = 2\,g(u_r)$ exatamente, em que $u_r =
        \omega_0^2 - 2\lambda^2$, e deduza que as duas
        frequências $\Omega_\pm$ em que $R = R_{\max}/\sqrt2$
        satisfazem $\Omega_+^2 - \Omega_-^2 =
        4\lambda\sqrt{\omega_0^2 - \lambda^2}$; conclua que, para
        amortecimento fraco, a largura de banda é $\Omega_+ - \Omega_-
        \approx 2\lambda$, isto é, $Q \approx
        \frac{\omega_0}{\Omega_+ - \Omega_-}$: picos de \omterm{ex:b1:diffeq:oscillation}{ressonância}
        agudos correspondem a sistemas de $Q$ alto.
  \item Retrato numérico para $\omega_0 = 1$, $\lambda = 0.05$
        ($Q = 10$), $A = 1$: calcule $\Omega_r$, $R_{\max}$, a
        resposta estática $R(0)$ e a largura de banda aproximada.
  \item Mostre que, se $2\lambda^2 \geq \omega_0^2$, então $R$ é
        estritamente decrescente em $\intoo0{+\infty}$: sistemas fortemente
        amortecidos não têm pico de \omterm{ex:b1:diffeq:oscillation}{ressonância} algum.
\end{enumerate}

\medskip
\noindent\textbf{Parte IV --- Sem amortecimento: batimentos e \omterm{ex:b1:diffeq:oscillation}{ressonância}.}
Aqui $\lambda = 0$.
\begin{enumerate}
  \item Para $\Omega \neq \omega_0$, encontre a solução geral de
        $x'' + \omega_0^2 x = A\cos(\Omega t)$.
  \item Resolva o \omterm{thm:b1:diffeq:voc}{problema de Cauchy} $x(0) = x'(0) = 0$ e transforme
        a resposta na forma de produto
        \[
        x(t) = \frac{2A}{\omega_0^2 - \Omega^2}\,
        \sin\Bigl(\frac{(\omega_0 - \Omega)t}2\Bigr)
        \sin\Bigl(\frac{(\omega_0 + \Omega)t}2\Bigr) .
        \]
  \item Para $\Omega$ próximo de $\omega_0$, leia o produto como uma
        oscilação rápida de frequência $\frac{\omega_0 + \Omega}2$
        modulada por um envelope lento de frequência
        $\frac{\abs{\omega_0 - \Omega}}2$: os \emph{batimentos}
        \index{batimentos}. Dê o período do envelope e a
        amplitude máxima, e note como ambos explodem quando $\Omega
        \to \omega_0$.
  \item Fixe $t$ e faça $\Omega \to \omega_0$ na fórmula da questão 19:
        mostre que o limite é
        \[
        x_\infty(t) = \frac{A\,t\,\sin(\omega_0 t)}{2\omega_0} ,
        \]
        e verifique diretamente que $x_\infty$ resolve a equação
        ressonante $x'' + \omega_0^2 x = A\cos(\omega_0 t)$ com
        $x(0) = x'(0) = 0$ (compare com o
        \cref{ex:b1:diffeq:oscillation}): a \omterm{ex:b1:diffeq:oscillation}{ressonância} é o limite
        de batimentos cada vez mais lentos e cada vez maiores.
  \item Contraste os dois destinos da \omterm{ex:b1:diffeq:oscillation}{ressonância}: crescimento linear
        $\frac{At}{2\omega_0}$ sem amortecimento, contra saturação
        em $R_{\max} \approx Q\,\frac{A}{\omega_0^2}$ com amortecimento
        fraco. Em uma frase: que mecanismo físico converte
        o primeiro no segundo?
\end{enumerate}

\medskip
\noindent\textbf{Parte V --- Balanço de energia e síntese.}
\begin{enumerate}
  \item No regime permanente da Parte II, calcule a média ao longo
        de um período $\frac{2\pi}\Omega$ de (a) a potência injetada
        pelo forçamento, $P_{\mathrm{in}}(t) = A\cos(\Omega t)
        \cdot x_p'(t)$, e (b) a potência dissipada pelo
        amortecimento, $P_{\mathrm{diss}}(t) = 2\lambda\,x_p'(t)^2$.
        Mostre que as duas médias valem $\lambda\,R^2\Omega^2$: o
        forçamento injeta exatamente o que o amortecimento queima ---
        eis por que o regime permanente é permanente.
  \item Onde exatamente o problema usou: (i) o teorema de estrutura
        \cref{thm:b1:diffeq:structure2}; (ii) o método da exponencial
        complexa; (iii) um estudo de função de variável real ao
        estilo do \cref{ch:b1:functions}? Uma frase para cada.
  \item Síntese: descreva o mapa completo de comportamentos de $(E_\Omega)$
        --- livre contra forçado, amortecido contra não amortecido, o
        papel de $Q$ como o único mostrador adimensional que ajusta
        altura do pico, largura de banda e tempo de vida do transitório
        --- e mencione onde a história continua: sistemas $2 \times 2$
        de primeira ordem (\cref{ch:b1:matrices} e o volume do segundo ano de
        graduação) e a decomposição de um forçamento periódico geral em
        senoides (séries de Fourier, no volume do terceiro ano de
        graduação), para as quais o caso senoidal deste problema é o
        tijolo fundamental.
\end{enumerate}
\end{problem}
